\documentclass[aspectratio=169,10pt]{beamer}

% =========================================================
% pdfLaTeX konfigūracija
% =========================================================

% Lietuviški simboliai
\usepackage[T1]{fontenc}
\usepackage[utf8]{inputenc}

% Kalbos
% Paskutinė nurodyta kalba yra pagrindinė
\usepackage[english,lithuanian]{babel}

% Noto Sans šriftas, suderinamas su pdfLaTeX
\usepackage[sfdefault]{noto}

% Matematika
\usepackage{amsmath,amssymb,cancel}

% =========================================================
% Beamer tema ir spalvų paletė
% =========================================================

\usetheme{Madrid}
\usecolortheme{whale}

\setbeamertemplate{navigation symbols}{}

% =========================================================
% Lietuviški aplinkų pavadinimai
% =========================================================

\deftranslation[to=lithuanian]{Definition}{Apibrėžimas}
\deftranslation[to=lithuanian]{Example}{Pavyzdys}
\deftranslation[to=lithuanian]{Remark}{Pastaba}
\deftranslation[to=lithuanian]{Proof}{Įrodymas}
\deftranslation[to=lithuanian]{Corollary}{Išvada}

% =========================================================
% Papildomos aplinkos
% =========================================================

\newenvironment{proposition}[1][]%
{\begin{block}{Teiginys: #1}}%
{\end{block}}

\newenvironment{property}[1][]%
{\begin{block}{Savybė: #1}}%
{\end{block}}

\newcommand{\vocab}[1]{\textbf{\color{structure.fg}#1}}

% =========================================================
% Dokumento informacija
% =========================================================

\title[Laipsnis su sveikuoju rodikliu]
{Laipsnis su sveikuoju rodikliu}

\subtitle{Teorija, pavyzdžiai ir praktinės užduotys}

\author[A. Berniukevičius]{Andrius Berniukevičius}

\date{\today}

% =========================================================
% Dokumentas
% =========================================================

\begin{document}

% ---------------------------------------------------------
% Titulinė skaidrė
% ---------------------------------------------------------

\begin{frame}
    \titlepage
\end{frame}

% ---------------------------------------------------------
% Nulinis laipsnis
% ---------------------------------------------------------

\begin{frame}{Nulinio laipsnio rodiklio atsiradimas}

    \begin{proposition}[Nulinis laipsnis]
        Jei $a \neq 0$, tai
        \[
        a^0 = 1.
        \]
    \end{proposition}

    \begin{proof}

        Pritaikykime laipsnių dalybos savybę

        \[
        \frac{a^m}{a^n}=a^{m-n},
        \]

        kai $m=n$:

    \end{proof}

\end{frame}

\begin{frame}{Nulinio laipsnio rodiklio atsiradimas}

    \begin{proof}

        \begin{enumerate}

            \item
            \textbf{Aritmetinė taisyklė:}
            bet kurį nelygų nuliui reiškinį padaliję iš savęs,
            gauname vienetą:

            \[
            \frac{a^n}{a^n}=1
            \]

            \item
            \textbf{Laipsnių dalybos taisyklė:}
            atimame rodiklius:

            \[
            \frac{a^n}{a^n}
            =
            a^{n-n}
            =
            a^0
            \]

        \end{enumerate}

        Kadangi pertvarkėme tą patį reiškinį,
        rezultatai yra lygūs:

        \[
        a^0=1
        \]

    \end{proof}

\end{frame}

% ---------------------------------------------------------
% Neigiamas rodiklis
% ---------------------------------------------------------

\begin{frame}{Laipsnis su neigiamu sveikuoju rodikliu}

    \begin{proposition}[Neigiamas rodiklis]

        Jei $a\neq0$ ir $k\in\mathbb{N}$, tai

        \[
        a^{-k}=\frac{1}{a^k}.
        \]

    \end{proposition}

    \begin{proof}

        Sudarykime trupmeną

        \[
        \frac{a^m}{a^{m+k}}
        \]

        (kur $m\in\mathbb{N}$)
        ir pertvarkykime dviem būdais:

    \end{proof}

\end{frame}

\begin{frame}{Laipsnis su neigiamu sveikuoju rodikliu}

    \begin{proof}

        \begin{enumerate}

            \item
            \textbf{Prastinimas:}

            išskaidome vardiklį

            \[
            a^{m+k}=a^m\cdot a^k
            \]

            ir suprastiname $a^m$:

            \[
            \frac{a^m}{a^{m+k}}
            =
            \frac{\cancelto{1}{a^m}}
            {\cancelto{1}{a^m}\cdot a^k}
            =
            \frac{1}{a^k}
            \]

            \item
            \textbf{Dalybos savybė:}
            atimame rodiklius:

            \[
            \frac{a^m}{a^{m+k}}
            =
            a^{m-(m+k)}
            =
            a^{m-m-k}
            =
            a^{-k}
            \]

        \end{enumerate}

    \end{proof}

\end{frame}

\begin{frame}{Laipsnis su neigiamu sveikuoju rodikliu}

    \begin{proof}

        Sulyginę gautus rezultatus, turime:

        \[
        a^{-k}=\frac{1}{a^k}.
        \]

    \end{proof}

\end{frame}

% ---------------------------------------------------------
% Trupmenos kėlimas neigiamu laipsniu
% ---------------------------------------------------------

\begin{frame}{Trupmenos kėlimas neigiamu laipsniu}

    \begin{block}{Išvada: Trupmenos apvertimas}

        Keliant trupmeną neigiamu laipsniu,
        ji yra apverčiama, o rodiklis tampa teigiamas:

        \[
        \left(\frac{a}{b}\right)^{-k}
        =
        \left(\frac{b}{a}\right)^k
        =
        \frac{b^k}{a^k}
        \qquad
        (a\neq0,\;b\neq0)
        \]

    \end{block}

    \begin{proof}

        Tarkime

        \[
        x=\frac{a}{b}
        \]

        ir pritaikykime formulę

        \[
        x^{-k}=\frac{1}{x^k}.
        \]

    \end{proof}

\end{frame}

\begin{frame}{Trupmenos kėlimas neigiamu laipsniu}

    \begin{proof}

        Tuomet

        \[
        \left(\frac{a}{b}\right)^{-k}
        =
        \frac{1}
        {\left(\frac{a}{b}\right)^k}
        =
        \frac{1}
        {\frac{a^k}{b^k}}
        =
        1:\frac{a^k}{b^k}
        =
        1\cdot\frac{b^k}{a^k}
        =
        \left(\frac{b}{a}\right)^k
        \]

    \end{proof}

    \begin{alertblock}{Taisyklė praktikai}

        Neigiamas rodiklis keičia dauginamojo vietą trupmenoje:
        iš skaitiklio keliauja į vardiklį
        (ir atvirkščiai) su teigiamu rodikliu.

    \end{alertblock}

\end{frame}

% ---------------------------------------------------------
% Apibrėžimas
% ---------------------------------------------------------

\begin{frame}{Laipsnio su sveikuoju rodikliu apibrėžimas}

    \begin{definition}[Laipsnis su sveikuoju rodikliu]

        \vocab{Laipsnis su sveikuoju rodikliu}
        yra reiškinys $a^n$, kur

        \[
        a\neq0,
        \qquad
        n\in\mathbb{Z}.
        \]

    \end{definition}

    \vspace{0.2cm}

    \begin{columns}[T]

        \begin{column}{0.32\textwidth}

            \begin{block}{1. Teigiamas ($n>0$)}

                \[
                a^n
                =
                \underbrace{
                a\cdot a\cdot\dots\cdot a
                }_{n\text{ kartų}}
                \]

                Įprasta $n$ vienodų dauginamųjų sandauga.

            \end{block}

        \end{column}

        \begin{column}{0.32\textwidth}

            \begin{block}{2. Nulinis ($n=0$)}

                \[
                a^0=1
                \]

                Bet koks nelygus nuliui skaičius,
                pakeltas $0$ laipsniu, lygus $1$.

            \end{block}

        \end{column}

        \begin{column}{0.32\textwidth}

            \begin{block}{3. Neigiamas ($n=-k$)}

                \[
                a^{-k}=\frac{1}{a^k}
                \]

                Reikšmė atvirkštinė laipsniui
                su teigiamu rodikliu.

            \end{block}

        \end{column}

    \end{columns}

\end{frame}

% ---------------------------------------------------------
% Lygtis ir trupmenos prastinimas
% ---------------------------------------------------------

\begin{frame}{Pavyzdžiai: Lygtis ir trupmenos prastinimas}

    \begin{columns}[T]

        \begin{column}{0.48\textwidth}

            \begin{example}[1]

                Išspręskite lygtį:

                \[
                \left(\frac{2}{5}\right)^x
                =
                6\frac{1}{4}
                \]

            \end{example}

            \begin{block}{Sprendimas}

                \begin{itemize}

                    \item
                    Mišrų skaičių verčiame trupmena:

                    \[
                    6\frac14
                    =
                    \frac{25}{4}
                    =
                    \left(\frac52\right)^2
                    \]

                    \item
                    Apverčiame pagrindą:

                    \[
                    \left(\frac52\right)^2
                    =
                    \left(\frac25\right)^{-2}
                    \]

                    \item
                    Sulyginame rodiklius:

                    \[
                    x=-2.
                    \]

                \end{itemize}

            \end{block}

        \end{column}

        \begin{column}{0.48\textwidth}

            \begin{example}[2]

                Supaprastinkite:

                \[
                \frac{28a^6b^4c^{-13}}
                {24a^{10}b^{-3}c}
                \]

            \end{example}

            \begin{block}{Sprendimas}

                \begin{align*}
                    &=
                    \frac{28}{24}
                    \cdot
                    a^{6-10}
                    \cdot
                    b^{4-(-3)}
                    \cdot
                    c^{-13-1}
                    \\[3pt]
                    &=
                    \frac76
                    \cdot
                    a^{-4}
                    \cdot
                    b^7
                    \cdot
                    c^{-14}
                    \\[3pt]
                    &=
                    \frac{7b^7}{6a^4c^{14}}
                \end{align*}

            \end{block}

        \end{column}

    \end{columns}

\end{frame}

% ---------------------------------------------------------
% Reiškinio pertvarkymas
% ---------------------------------------------------------

\begin{frame}{Pavyzdžiai: Reiškinio pertvarkymas}

    \small

    \begin{example}[3]

        Supaprastinkite:

        \[
        \left(
        \frac{3}{10}x^{-4}yz^2
        \right)^{-4}
        \]

        \[
        \left(
        \frac{3yz^2}{10x^4}
        \right)^{-4}
        =
        \left(
        \frac{10x^4}{3yz^2}
        \right)^4
        =
        \frac{
        10^4\cdot(x^4)^4
        }{
        3^4\cdot y^4\cdot(z^2)^4
        }
        =
        \frac{10000x^{16}}
        {81y^4z^8}
        \]

    \end{example}

    \begin{example}[4]

        Supaprastinkite:

        \[
        \left(
        \frac{6x^{-2}y^3}{5ab^4}
        \right)^{-2}
        :
        \frac{5a^{-3}b}{72xy^4}
        \]

        \[
        \left(
        \frac{5ab^4x^2}{6y^3}
        \right)^2
        \cdot
        \frac{72a^3xy^4}{5b}
        =
        \frac{25a^2b^8x^4}{36y^6}
        \cdot
        \frac{72a^3xy^4}{5b}
        =
        \frac{10a^5b^7x^5}{y^2}
        \]

    \end{example}

\end{frame}

% ---------------------------------------------------------
% Uždaviniai 1--12
% ---------------------------------------------------------

\begin{frame}{Uždaviniai (1 dalis: 1--12)}

    \small

    Atlikite veiksmus
    (atsakymus pateikite tik su teigiamais rodikliais):

    \vspace{0.2cm}

    \begin{columns}[T]

        \begin{column}{0.5\textwidth}

            \begin{enumerate}

                \item
                $6^x=\frac{1}{216}$

                \item
                $\left(\frac34\right)^x
                =
                2\frac{10}{27}$

                \item
                $0{,}2^x=125$

                \item
                $\left(
                \frac34a^3b^{-2}
                \right)^{-2}$

                \item
                $\left(
                \frac{2x^3}{3y^4}
                \right)^{-4}$

                \item
                $\left(
                \frac{2m^3}{5n^7}
                \right)^{-2}
                \cdot
                \frac{4m^3}{125n^{10}}$

            \end{enumerate}

        \end{column}

        \begin{column}{0.5\textwidth}

            \begin{enumerate}

                \setcounter{enumi}{6}

                \item
                $2{,}5^x=\frac{8}{125}$

                \item
                $5^x=0{,}0016$

                \item
                $\left(
                0{,}1x^5y^{-4}z
                \right)^{-3}$

                \item
                $\frac{
                2x^3y^{-6}z
                }{
                12x^4y^6z^{-3}
                }$

                \item
                $\frac{
                36a^5bc^{-10}d^3
                }{
                24ab^{-1}c^{-3}d^6
                }$

                \item
                $\frac{
                45mn^{10}p^{-3}
                }{
                25m^{-3}n^0p^4
                }$

            \end{enumerate}

        \end{column}

    \end{columns}

\end{frame}

% ---------------------------------------------------------
% Uždaviniai 13--23
% ---------------------------------------------------------

\begin{frame}{Uždaviniai (2 dalis: 13--23)}

    \begin{columns}[T]

        \begin{column}{0.5\textwidth}

            \begin{enumerate}

                \setcounter{enumi}{12}

                \item
                $\left(
                \frac{3ab^4}{5xy}
                \right)^{-1}
                \cdot
                \frac{9x^{-1}y^2}{10ab^4}$

                \item
                $4^x=\frac{1}{64}$

                \item
                $\left(
                \frac23
                \right)^x
                =
                \frac{27}{8}$

                \item
                $0{,}25^x=64$

                \item
                $\left(
                \frac52c^{-2}d^3
                \right)^{-2}$

                \item
                $\left(
                \frac{3p^2}{4q^3}
                \right)^{-3}$

            \end{enumerate}

        \end{column}

        \begin{column}{0.5\textwidth}

            \begin{enumerate}

                \setcounter{enumi}{18}

                \item
                $\left(
                \frac{7r^{-3}s^2}{2t}
                \right)^{-2}
                \cdot
                \frac{49r^2}{4s^4}$

                \item
                $0{,}4^x=\frac{125}{8}$

                \item
                $\left(
                \frac13
                \right)^x
                =
                81$

                \item
                $\left(
                0{,}2u^{-3}v^2w^{-1}
                \right)^{-2}$

                \item
                $\left(
                \frac{4a^{-2}b^3}{9c}
                \right)^{-1}
                :
                \frac{27ac^2}{8b^5}$

            \end{enumerate}

        \end{column}

    \end{columns}

\end{frame}

% ---------------------------------------------------------
% Atsakymai
% ---------------------------------------------------------

\begin{frame}{Atsakymai (1--23)}

    \begin{columns}[T]

        \begin{column}{0.33\textwidth}

            \small

            \textbf{(1)} $x=-3$ \\[3pt]

            \textbf{(2)} $x=-3$ \\[3pt]

            \textbf{(3)} $x=-3$ \\[3pt]

            \textbf{(4)}
            $\dfrac{16b^4}{9a^6}$ \\[5pt]

            \textbf{(5)}
            $\dfrac{81y^{16}}{16x^{12}}$ \\[5pt]

            \textbf{(6)}
            $\dfrac{n^4}{5m^3}$ \\[5pt]

            \textbf{(7)} $x=-3$ \\[3pt]

            \textbf{(8)} $x=-4$

        \end{column}

        \begin{column}{0.33\textwidth}

            \small

            \textbf{(9)}
            $\dfrac{1000y^{12}}{x^{15}z^3}$ \\[5pt]

            \textbf{(10)}
            $\dfrac{z^4}{6xy^{12}}$ \\[5pt]

            \textbf{(11)}
            $\dfrac{3a^4b^2}{2c^7d^3}$ \\[5pt]

            \textbf{(12)}
            $\dfrac{9m^4n^{10}}{5p^7}$ \\[5pt]

            \textbf{(13)}
            $\dfrac{3y^3}{2a^2b^8}$ \\[5pt]

            \textbf{(14)} $x=-3$ \\[3pt]

            \textbf{(15)} $x=-3$ \\[3pt]

            \textbf{(16)} $x=-3$

        \end{column}

        \begin{column}{0.33\textwidth}

            \small

            \textbf{(17)}
            $\dfrac{4c^4}{25d^6}$ \\[5pt]

            \textbf{(18)}
            $\dfrac{64q^9}{27p^6}$ \\[5pt]

            \textbf{(19)}
            $\dfrac{r^8t^2}{s^8}$ \\[5pt]

            \textbf{(20)} $x=-3$ \\[3pt]

            \textbf{(21)} $x=-4$ \\[3pt]

            \textbf{(22)}
            $\dfrac{25u^6w^2}{v^4}$ \\[5pt]

            \textbf{(23)}
            $\dfrac{2ab^2}{3c}$

        \end{column}

    \end{columns}

\end{frame}

\end{document}